% =============================================================
%  06_macro.tex
%  Capitolo 6 --- Macro
% =============================================================

\chapter{Macro}

\begin{intuizione}
Le funzioni trasformano \emph{valori}.
Le macro trasformano \emph{codice}.
Una macro riceve espressioni non valutate e le riscrive
in altre espressioni, prima che il compilatore le veda.
Questo è il meccanismo con cui Lisp si estende: puoi aggiungere
nuove strutture di controllo al linguaggio, indistinguibili
da quelle native.
\end{intuizione}

\vspace{4pt}

% ================================================================
\section{La differenza tra funzione e macro}
% ================================================================

\begin{concetto}{Funzione vs macro: quando vengono valutati gli argomenti}
\keyline{}{Funzione: gli argomenti sono valutati \emph{prima} della chiamata.}
\keyline{}{Macro: gli argomenti arrivano come \emph{codice grezzo}, non valutati.}

\boxsep
\detail{Questa differenza è fondamentale.
Una funzione non può implementare \fn{if}: se tutti gli argomenti
fossero valutati prima della chiamata, entrambi i rami verrebbero
eseguiti sempre. \fn{if} è una \emph{forma speciale} (special operator):
possiede regole di valutazione proprie che controllano \emph{quale}
ramo valutare. Forme come \fn{when}, \fn{unless}, \fn{and}, \fn{or}
sono invece macro costruite sopra \fn{if}.}
\end{concetto}

\begin{tcolorbox}[
  enhanced,
  colback=black!2, colframe=black!20,
  arc=3pt, boxrule=0.4pt,
  left=12pt, right=12pt, top=9pt, bottom=9pt,
  before skip=8pt, after skip=10pt
]
\small\color{black!65}

\textbf{Perché le funzioni non bastano.}

\medskip
Supponi di voler implementare \fn{mio-and} come funzione:

\begin{tcolorbox}[enhanced, colback=clBg, colframe=gray!40,
  arc=2pt, boxrule=0.4pt, left=8pt, right=6pt, top=4pt, bottom=4pt,
  before skip=4pt, after skip=4pt]
\begin{lstlisting}
;; SBAGLIATO: non funziona come and
(defun mio-and-fn (a b)
  (if a b nil))

;; Il problema:
(mio-and-fn t (/ 1 0))   ; ERRORE: divide by zero!
;; Perche'? Prima di chiamare mio-and-fn, Lisp valuta TUTTI
;; gli argomenti. (/ 1 0) viene valutato anche se a e' NIL.

;; and nativo usa short-circuit: non valuta b se a e' falso
(and nil (/ 1 0))   ; => NIL  (nessun errore)
\end{lstlisting}
\end{tcolorbox}

\medskip
La valutazione cortocircuito di \fn{and} non si può implementare
con una funzione. Si può implementare con una macro, perché la macro
riceve le espressioni non valutate e decide lei come (e se) valutarle.

\medskip
\textbf{Questo è il potere delle macro}: non si limitano a calcolare
con valori, possono \emph{ristrutturare} il codice prima
che venga eseguito.
In Java o Python non esiste nulla di equivalente.
Le annotazioni, i decorator, il preprocessore C --- sono tutti
meccanismi più limitati che operano in fasi separate.
Le macro Lisp lavorano direttamente sulle forme del programma,
rappresentate come dati (liste, simboli, atomi), con tutto il
linguaggio disponibile per manipolarle.
\end{tcolorbox}

% ================================================================
\section{Come funziona una macro}
% ================================================================

\begin{concetto}{Il ciclo di vita di una macro}
\keyline{}{1. Il compilatore incontra \fn{(nome-macro arg1 arg2 ...)}.}
\keyline{}{2. Chiama il \emph{trasformatore} della macro con gli argomenti\\
come liste simboliche (codice grezzo).}
\keyline{}{3. Il trasformatore restituisce una nuova espressione S\\
(fase di \emph{macro expansion}).}
\keyline{}{4. La forma risultante viene compilata o valutata\\
secondo il contesto.}

\boxsep
\detail{Il trasformatore è codice Common Lisp ordinario.
Può usare tutto il linguaggio: variabili, funzioni, ricorsione.
L'espansione avviene a \emph{compile time} (o prima della valutazione
nell'interprete), non a runtime.}
\end{concetto}

\begin{esempio}[Prima macro: \fn{quando} come \fn{when}]
\begin{lstlisting}
;; defmacro: come defun, ma riceve codice grezzo
(defmacro quando (test &body corpo)
  ;; test e corpo sono simboli/liste, non valori
  ;; questo codice gira a compile-time
  `(if ,test (progn ,@corpo) nil))

;; Uso
(quando (> x 0)
  (print "positivo")
  (print "fatto"))

;; Equivale a:
(if (> x 0)
    (progn (print "positivo")
           (print "fatto"))
    nil)
\end{lstlisting}
\end{esempio}

\begin{attenzione}
\textbf{\fn{macroexpand-1}: il tuo migliore amico.}
Prima di usare una macro in produzione, verifica sempre la sua
espansione:
\begin{lstlisting}
(macroexpand-1 '(quando (> x 0) (print "positivo")))
; => (IF (> X 0) (PROGN (PRINT "positivo")) NIL)  T
\end{lstlisting}
\fn{macroexpand-1} espande un solo livello.
\fn{macroexpand} espande tutto ricorsivamente.
\end{attenzione}

% ================================================================
\section{Backquote: costruire codice programmaticamente}
% ================================================================

\begin{concetto}{Backquote, comma e comma-at}
\keyline{\texttt{`(a b c)}}{come \fn{'(a b c)} ma con interpolazione.}
\keyline{\texttt{,expr}}{valuta \fn{expr} e inserisce il risultato.}
\keyline{\texttt{,@lista}}{valuta \fn{lista} e la \emph{espande} (splice).}
\end{concetto}

\begin{tcolorbox}[
  enhanced,
  colback=black!2, colframe=black!20,
  arc=3pt, boxrule=0.4pt,
  left=12pt, right=12pt, top=9pt, bottom=9pt,
  before skip=8pt, after skip=10pt
]
\small\color{black!65}

\textbf{Il backquote è il template engine di Lisp.}

\medskip
Senza backquote, costruire una lista che rappresenta codice
è verboso e error-prone:

\begin{tcolorbox}[enhanced, colback=clBg, colframe=gray!40,
  arc=2pt, boxrule=0.4pt, left=8pt, right=6pt, top=4pt, bottom=4pt,
  before skip=4pt, after skip=4pt]
\begin{lstlisting}
;; Senza backquote: costruzione manuale
(list 'if test (cons 'progn corpo) nil)

;; Con backquote: template leggibile
`(if ,test (progn ,@corpo) nil)
\end{lstlisting}
\end{tcolorbox}

\medskip
Il \textbf{backquote} (\texttt{`}) avvia un template.
All'interno, tutto è quotato tranne ciò che è preceduto da:
\begin{itemize}
  \item \textbf{\texttt{,}} (comma) --- inserisce il valore di
        un'espressione
  \item \textbf{\texttt{,@}} (comma-at) --- inserisce gli elementi
        di una lista (splice)
\end{itemize}

\begin{tcolorbox}[enhanced, colback=clBg, colframe=gray!40,
  arc=2pt, boxrule=0.4pt, left=8pt, right=6pt, top=4pt, bottom=4pt,
  before skip=4pt, after skip=4pt]
\begin{lstlisting}
(let ((x 42)
      (lst '(1 2 3)))
  `(a ,x c)       ; => (A 42 C)      , inserisce il valore
  `(a ,@lst c)    ; => (A 1 2 3 C)   ,@ splica la lista
  `(a ,(+ 1 2))   ; => (A 3)         , valuta l'espressione
  `(a 'b)         ; => (A (QUOTE B)) senza , resta quotato
)
\end{lstlisting}
\end{tcolorbox}

\medskip
La differenza tra \texttt{,x} e \texttt{,@x}:
\texttt{,x} inserisce \texttt{x} come \emph{un elemento};
\texttt{,@x} inserisce gli \emph{elementi di x}.

\begin{tcolorbox}[enhanced, colback=clBg, colframe=gray!40,
  arc=2pt, boxrule=0.4pt, left=8pt, right=6pt, top=4pt, bottom=4pt,
  before skip=4pt, after skip=4pt]
\begin{lstlisting}
(let ((corpo '((print "a") (print "b"))))
  `(progn ,corpo)    ; => (PROGN ((PRINT "a") (PRINT "b")))
  `(progn ,@corpo))  ; => (PROGN (PRINT "a") (PRINT "b"))
;; ,@corpo inserisce le due espressioni come elementi separati
\end{lstlisting}
\end{tcolorbox}
\end{tcolorbox}

% ================================================================
\section{Macro pratiche}
% ================================================================

\argomento{\fn{\&body}: raccogliere il corpo}

\begin{concetto}{\fnt{\&body}}
\keyline{}{Come \fn{\&rest} ma segnala a editor e strumenti\\
che questi argomenti sono il \emph{corpo} della macro.}
\detail{La differenza è solo di indentazione automatica:
gli editor Lisp indentano correttamente le forme con \fn{\&body}.}
\end{concetto}

\begin{esempio}[Macro \fn{mio-while}]
\begin{lstlisting}
(defmacro mio-while (test &body corpo)
  `(loop
     (unless ,test (return))
     ,@corpo))

;; Uso
(let ((i 0))
  (mio-while (< i 5)
    (print i)
    (setf i (1+ i))))
; stampa: 0 1 2 3 4

;; Espansione
(macroexpand-1 '(mio-while (< i 5) (print i)))
; => (LOOP (UNLESS (< I 5) (RETURN)) (PRINT I))
\end{lstlisting}
\end{esempio}

\begin{workbook}
\textbf{Esercizi 121--123}
(\fn{quando}, \fn{mio-while}, \fn{mio-for})\\[4pt]
Questi tre esercizi seguono una progressione precisa.
Dopo ciascuno usa \fn{macroexpand-1} per verificare l'espansione ---
è l'unico modo per essere sicuri che la macro faccia quello che pensi.\\[4pt]
L'esercizio~123 (\fn{mio-for}) introduce \fn{gensym}: leggilo con
attenzione anche se non lo implementi, perché il problema della
cattura di variabile è reale e importante.
\end{workbook}

% ================================================================
\section{Il problema della cattura di variabile}
% ================================================================

\begin{concetto}{Cattura di variabile --- il bug silenzioso delle macro}
\keyline{}{Se la macro usa nomi di variabile ``comuni'', può\\
\emph{catturare} variabili del codice chiamante per errore.}
\end{concetto}

\begin{tcolorbox}[
  enhanced,
  colback=black!2, colframe=black!20,
  arc=3pt, boxrule=0.4pt,
  left=12pt, right=12pt, top=9pt, bottom=9pt,
  before skip=8pt, after skip=10pt
]
\small\color{black!65}

\textbf{Il problema.}

\medskip
Considera questa macro \fn{swap!} che scambia due variabili:

\begin{tcolorbox}[enhanced, colback=clBg, colframe=gray!40,
  arc=2pt, boxrule=0.4pt, left=8pt, right=6pt, top=4pt, bottom=4pt,
  before skip=4pt, after skip=4pt]
\begin{lstlisting}
;; VERSIONE CON BUG
(defmacro swap!-bug (a b)
  `(let ((tmp ,a))      ; tmp potrebbe essere catturato!
     (setf ,a ,b)
     (setf ,b tmp)))

;; Funziona nella maggior parte dei casi...
(let ((x 1) (y 2))
  (swap!-bug x y)
  (list x y))         ; => (2 1)  ok

;; ...ma non se il chiamante usa una variabile chiamata tmp!
(let ((tmp 10) (y 20))
  (swap!-bug tmp y)
  (list tmp y))       ; => SBAGLIATO: tmp viene catturato
\end{lstlisting}
\end{tcolorbox}

\medskip
\textbf{La soluzione: \fn{gensym}.}

\fn{gensym} genera un simbolo unico garantito mai usato altrove.

\begin{tcolorbox}[enhanced, colback=clBg, colframe=gray!40,
  arc=2pt, boxrule=0.4pt, left=8pt, right=6pt, top=4pt, bottom=4pt,
  before skip=4pt, after skip=4pt]
\begin{lstlisting}
;; VERSIONE CORRETTA con gensym
(defmacro swap! (a b)
  (let ((tmp (gensym)))   ; gensym gira a compile-time
    `(let ((,tmp ,a))     ; ,tmp inserisce il simbolo unico
       (setf ,a ,b)
       (setf ,b ,tmp))))

;; Ora e' sicuro
(let ((tmp 10) (y 20))
  (swap! tmp y)
  (list tmp y))   ; => (20 10)  corretto

;; Espansione: il nome e' qualcosa come G1234
(macroexpand-1 '(swap! x y))
; => (LET ((#:G1234 X)) (SETF X Y) (SETF Y #:G1234))
\end{lstlisting}
\end{tcolorbox}

\medskip
\textbf{Regola pratica}: ogni variabile temporanea introdotta
da una macro deve usare \fn{gensym}.
Se la macro non introduce variabili temporanee (come \fn{quando}
o \fn{mio-while}), il problema non si pone.
\end{tcolorbox}

% ================================================================
\section{Quando usare le macro}
% ================================================================

\begin{concetto}{Macro: sì o no?}
\keyline{}{Usa una macro \emph{solo} quando una funzione non è sufficiente.}

\boxsep
\detail{La tentazione di usare le macro è forte in Lisp,
ma di solito una funzione o una higher-order function sono
la scelta giusta. Le macro aggiungono complessità e
rendono il codice più difficile da analizzare.}
\end{concetto}

\begin{tcolorbox}[
  enhanced,
  colback=black!2, colframe=black!20,
  arc=3pt, boxrule=0.4pt,
  left=12pt, right=12pt, top=9pt, bottom=9pt,
  before skip=8pt, after skip=10pt
]
\small\color{black!65}

\textbf{Usa una macro quando:}
\begin{itemize}
  \item Hai bisogno di \textbf{valutazione controllata}: alcuni argomenti
        non devono essere valutati, o devono essere valutati
        in un certo ordine (\fn{if}, \fn{and}, \fn{or},
        \fn{when}, \fn{while}).
  \item Vuoi introduire \textbf{nuova sintassi} o nuovi costrutti
        di controllo del flusso.
  \item Vuoi evitare la ripetizione di \textbf{boilerplate} strutturale
        che una funzione non può eliminare (es.\ \fn{with-open-file}
        gestisce automaticamente la chiusura del file).
\end{itemize}

\medskip
\textbf{Non usare una macro quando:}
\begin{itemize}
  \item Una funzione di ordine superiore risolve il problema
        (preferisci \fn{mapcar \#'f lst} a una macro che fa la stessa cosa).
  \item Il ``vantaggio'' è solo sintattico e non vale la complessità aggiunta.
  \item Il codice diventa difficile da testare o da \fn{macroexpand}.
\end{itemize}

\medskip
\textbf{Esempi di macro nella libreria standard} che mostrano
quando le macro sono giustificate:
\begin{itemize}
  \item \fn{with-open-file} --- garantisce la chiusura del file
        anche in caso di errore (il corpo è valutato in un contesto
        protetto con cleanup automatico)
  \item \fn{loop} --- mini-linguaggio per iterazione
  \item \fn{setf} --- generalizzazione dell'assegnazione a qualsiasi
        ``posto'' (lista, array, slot di struttura...)
  \item \fn{defun}, \fn{defstruct}, \fn{defclass} --- definizioni
        che creano binding globali e metadata
\end{itemize}
\end{tcolorbox}

\begin{esempio}[\fn{with-open-file}: macro nella pratica]
\begin{lstlisting}
;; with-open-file e' una macro: garantisce la chiusura
(with-open-file (stream "dati.txt" :direction :input)
  (read-line stream))
;; Anche se read-line segnala un errore, il file viene chiuso.

;; Espansione semplificata:
;; (let ((stream (open "dati.txt" :direction :input)))
;;   (unwind-protect
;;     (read-line stream)
;;     (close stream)))
;;
;; unwind-protect esegue il cleanup anche in caso di eccezione.
;; Senza macro, il chiamante dovrebbe scrivere tutto questo ogni volta.
\end{lstlisting}
\end{esempio}

% ================================================================
\section*{Riepilogo del capitolo}
% ================================================================

\begin{tcolorbox}[
  enhanced,
  colback=csBlue!5, colframe=csBlue!40,
  arc=4pt, boxrule=0.6pt,
  left=14pt, right=14pt, top=12pt, bottom=12pt,
  before skip=16pt, after skip=16pt
]
\textbf{Quello che hai imparato:}

\medskip
\begin{itemize}
  \item \textbf{Macro vs funzione}: le funzioni ricevono valori,
        le macro ricevono codice grezzo non valutato.
        Le macro vengono espanse a compile-time, le funzioni
        girano a runtime.
  \item \textbf{\fn{defmacro}}: come \fn{defun} ma con parametri
        che sono S-expression, non valori.
        \fn{\&body} per il corpo della macro.
  \item \textbf{Backquote}: template per costruire codice.
        \texttt{,expr} inserisce un valore,
        \texttt{,@lista} espande una lista.
  \item \textbf{\fn{macroexpand-1}}: verifica sempre l'espansione.
        È lo strumento di debug fondamentale per le macro.
  \item \textbf{Cattura di variabile}: usa \fn{gensym} per ogni
        variabile temporanea introdotta dalla macro.
  \item \textbf{Quando usarle}: valutazione controllata, nuova
        sintassi, eliminazione di boilerplate strutturale.
        Non usarle se una funzione o una HOF bastano.
\end{itemize}

\medskip
\textbf{Prossimo:} Capitolo~7 --- \emph{Backtracking e algoritmi}.
Spazi di ricerca, sottoinsiemi, permutazioni, ordinamento,
problemi classici.
\end{tcolorbox}

\begin{workbook}
\textbf{Prima di andare al Capitolo~7:}\\[4pt]
Completa gli esercizi \textbf{121--123} (capitolo \emph{Macro}).\\[4pt]
Sono solo tre esercizi, ma richiedono un cambio di prospettiva:
pensare al codice come dato.
Usa \fn{macroexpand-1} dopo ogni esercizio, anche se la macro
sembra funzionare: assicurati di capire esattamente cosa genera.
\end{workbook}
